\documentclass[final]{beamer}

% 42 inches wide x 36 inches tall (106.68 cm x 91.44 cm), landscape.
\usepackage[size=custom,width=106.68,height=91.44,scale=0.92]{beamerposter}
\usepackage[T1]{fontenc}
\usepackage[utf8]{inputenc}
\usepackage{lmodern}
\usepackage{amsmath,amssymb,mathtools}
\usepackage{booktabs,array}
\usepackage{xcolor}
\usepackage{tikz}
\usetikzlibrary{positioning,arrows.meta,calc,fit,shapes.geometric}
\usepackage{graphicx}
\usepackage{ragged2e}
\usepackage{url}
\usepackage{hyperref}

% -------------------------------------------------------------------------
% Colors and typography
% -------------------------------------------------------------------------
\definecolor{WashURed}{HTML}{BA0C2F}
\definecolor{WashUGreen}{HTML}{215732}
\definecolor{PosterDark}{HTML}{222222}
\definecolor{PosterGray}{HTML}{666666}
\definecolor{PosterLight}{HTML}{F6F6F6}
\definecolor{PosterPaleGreen}{HTML}{EAF2EC}
\definecolor{PosterPaleRed}{HTML}{F8E9ED}

\setbeamercolor{headline}{fg=white,bg=WashURed}
\setbeamercolor{title in headline}{fg=white}
\setbeamercolor{author in headline}{fg=white}
\setbeamercolor{institute in headline}{fg=white}
\setbeamercolor{block title}{fg=white,bg=WashUGreen}
\setbeamercolor{block body}{fg=PosterDark,bg=white}
\setbeamercolor{block alerted title}{fg=white,bg=WashURed}
\setbeamercolor{block alerted body}{fg=PosterDark,bg=PosterLight}
\setbeamercolor{item}{fg=WashURed}

\setbeamertemplate{navigation symbols}{}
\setbeamertemplate{itemize item}{\color{WashURed}$\blacktriangleright$}
\setbeamertemplate{itemize subitem}{\color{WashUGreen}$\bullet$}
\setbeamertemplate{caption}{\raggedright\insertcaption\par}

\setbeamerfont{title in headline}{size=\fontsize{66}{72}\selectfont,series=\bfseries}
\setbeamerfont{author in headline}{size=\fontsize{31}{37}\selectfont}
\setbeamerfont{institute in headline}{size=\fontsize{31}{37}\selectfont}
\setbeamerfont{block title}{size=\fontsize{35}{41}\selectfont,series=\bfseries}
\setbeamerfont{block body}{size=\fontsize{27}{33}\selectfont}
\setbeamerfont{caption}{size=\fontsize{18}{22}\selectfont}

% -------------------------------------------------------------------------
% Layout helpers
% -------------------------------------------------------------------------
\newlength{\sepwidth}
\newlength{\colwidth}
\setlength{\sepwidth}{0.0175\paperwidth}
\setlength{\colwidth}{0.31\paperwidth}
\newcommand{\separatorcolumn}{\begin{column}{\sepwidth}\end{column}}

\newcommand{\keybox}[1]{%
\begin{beamercolorbox}[rounded=true,shadow=false,sep=0.65em,wd=\linewidth]{block alerted body}
\centering #1
\end{beamercolorbox}}

\newcommand{\assumptionbox}[2]{%
\begin{beamercolorbox}[rounded=true,shadow=false,sep=0.6em,wd=\linewidth]{block alerted body}
\textbf{#1}\par\smallskip #2
\end{beamercolorbox}}

\newcommand{\threatbox}[2]{%
\begin{beamercolorbox}[rounded=true,shadow=false,sep=1.0em,wd=\linewidth]{block alerted body}
\centering
\parbox[t][5.2cm]{\dimexpr\linewidth-2.0em\relax}{%
  \justifying
  \textbf{#1}\par\vspace{0.5em} #2
}%
\end{beamercolorbox}}


\newcommand{\figureplaceholder}[2]{%
\begin{center}
\begin{tikzpicture}
  \draw[very thick,dashed,color=PosterGray,rounded corners=8pt] (0,0) rectangle (#1,#2);
  \node[align=center,color=PosterGray] at (#1/2,#2/2)
    {\Large Figure placeholder\\[0.3em]\normalsize Replace with \texttt{\textbackslash includegraphics}};
\end{tikzpicture}
\end{center}}

\newcommand{\maybeincludegraphics}[2]{%
\IfFileExists{#2}{\includegraphics[#1]{#2}}{\figureplaceholder{13.5}{5.7}}}

\title{Causal Inference with Ordinal Outcomes:\ Flexible Estimation Using Normalizing Flows}
\author{Chanhyuk Park}
\institute{Washington University in St. Louis}

\begin{document}
\begin{frame}[t]

% -------------------------------------------------------------------------
% Header
% -------------------------------------------------------------------------
\begin{beamercolorbox}[wd=\paperwidth,sep=0.7cm,leftskip=1.4cm,rightskip=1.4cm]{headline}
  \centering
  \usebeamerfont{title in headline}\vspace{1cm}\par\inserttitle\par\vspace{0.18cm}
  \usebeamerfont{author in headline}\insertauthor\par\vspace{0.03cm}
  \usebeamerfont{institute in headline}\insertinstitute\vspace{0.25cm}
\end{beamercolorbox}

\vspace{0.42cm}
\begin{columns}[t,totalwidth=0.995\paperwidth]
\separatorcolumn

% =========================================================================
% COLUMN 1: IDENTIFICATION
% =========================================================================
\begin{column}{\colwidth}

\begin{block}{Latent Traits and Ordinal Measurement}
\justifying
\begin{itemize}
  \item Political science wants to evaluate causal effects on continuous, latent traits $Y_i^*$ (e.g., support for policies, trust in institutions) {\footnotesize (McKelvey and Zavoina, 1975; Canes-Wrone and De Marchi, 2002; Kriner and Schwartz, 2009; Tomz and Weeks, 2020; Alt and Iversen, 2017; Magni, 2021)}.
  \item However, we routinely observe ordinal survey indicators $Y_i \in \{1, \dots, J\}$ (e.g., a Likert scale).
\end{itemize}

\vspace{0.3cm}
\textbf{Model and Variables:}
Let $D_i \in \{0, 1\}$ be a binary treatment indicator, $X_i$ be baseline covariates, $\epsilon_i$ be a latent error term with CDF $F(\cdot)$, and $\alpha$ be thresholds in the latent space. We model the continuous latent trait $Y_i^*$ using the systematic latent index $m(X_i, D_i)$ (which simplifies to $\tau D_i + X_i^\top \beta$ under a linear model with latent treatment effect $\tau$):
\[
Y_i^* = m(X_i, D_i) + \epsilon_i
\]
The observed categories arise from a threshold-crossing model partitioning the latent scale:
\[
Y_i=j \iff \alpha_{j-1}<Y_i^*\le \alpha_j, \qquad j=1,\dots,J
\]

  \vspace{0.5cm}
  % latent {{{
\begin{center}
\begin{tikzpicture}[x=2cm,y=1.5cm,>=Latex]
  % Draw x-axis
  \draw[->,thick] (-4.8,0)--(4.8,0) node[right] {$Y_i^*$};

  % Draw thresholds (2 thresholds for 3 categories)
  \draw[thick] (-2,-.12)--(-2,.12) node[below=3pt] {$\alpha_1$};
  \draw[thick] (2,-.12)--(2,.12) node[below=3pt] {$\alpha_2$};

  % Category labels at the top
  \node[below=18pt,font=\small,color=PosterDark] at (-3.25,0) {Category 1};
  \node[below=18pt,font=\small,color=PosterDark] at (0,0) {Category 2};
  \node[below=18pt,font=\small,color=PosterDark] at (3.25,0) {Category 3};

  % Control curve (wider)
  \draw[thick,WashUGreen] plot[smooth] coordinates {
    (-4.50,0.04) (-3.86,0.09) (-3.21,0.19) (-2.57,0.34) (-1.93,0.53) (-1.29,0.73) (-0.64,0.88) (0.00,0.94) (0.64,0.88) (1.29,0.73) (1.93,0.53) (2.57,0.34) (3.21,0.19) (3.86,0.09) (4.50,0.04)
  };

  % Treatment curve (narrower, higher peak)
  \draw[thick,WashURed,dashed] plot[smooth] coordinates {
    (-4.50,0.00) (-3.86,0.00) (-3.21,0.00) (-2.57,0.03) (-1.93,0.13) (-1.29,0.44) (-0.64,0.97) (0.00,1.40) (0.64,1.34) (1.29,0.84) (1.93,0.35) (2.57,0.10) (3.21,0.02) (3.86,0.00) (4.50,0.00)
  };

  % Text labels for curves
  \node[WashUGreen,above] at (-2.2,0.65) {Control};
  \node[WashURed,above] at (2,1.0) {Treatment};
\end{tikzpicture}
\end{center}
% }}}
\end{block}


\begin{block}{Why Cardinalization Fails (Toy Example: The Sign Flips)}
\justifying
A common approach is to assign arbitrary numerical scores $c(j)$ to the categories. The cardinalized treatment effect is then the difference in expected values:
\[
\tau_c = \mathbb E[c(Y_i(1)) - c(Y_i(0))] = \sum_{j=1}^J c(j) \left[ \Pr(Y_i(1) = j) - \Pr(Y_i(0) = j) \right].
\]

Suppose treatment redistributes the probability distribution across three ordered categories from Control to Treatment as follows:
\[
\Pr(Y_i(0) = j) = (0.30, \; 0.40, \; 0.30) \quad \longrightarrow \quad \Pr(Y_i(1) = j) = (0.15, \; 0.65, \; 0.20).
\]

Then two strictly monotone codings starting from 1 yield opposite signs:
\[
\begin{aligned}
 c_A=(1,2,3):\quad &\tau_{c_A}= 1 \cdot (-0.15) + 2 \cdot 0.25 + 3 \cdot (-0.10) =\mathbf{+0.05},\\
 c_B=(1,2,6):\quad &\tau_{c_B}= 1 \cdot (-0.15) + 2 \cdot 0.25 + 6 \cdot (-0.10) =\mathbf{-0.25}.
\end{aligned}
\]

\begin{itemize}\setlength\itemsep{0.25em}
  \item Treating a Likert scale as continuous is equivalent to imposing an arbitrary linear cardinalization.
  \item As shown by Liddell and Kruschke (2018) and Bond and Lang (2019), the estimated ATE based on such cardinalizations can introduce bias and even flip estimated signs. 
\end{itemize}
\end{block}

\vspace{-0.5cm}

\begin{block}{Failure of Latent Scale Identification}
\justifying
\begin{itemize}
  \item Because the location and scale of $Y_i^*$ are unobserved, we can apply any positive affine transformation that yields same observed ordinal distribution ($\tilde{Y}_i^* = a + c Y_i^*$ and $\tilde{\alpha}_j = a + c \alpha_j$ with $c > 0$):
\end{itemize}

  \vspace{-0.7cm}

\[
\Pr(\tilde{\alpha}_{j-1} < \tilde{Y}_i^* \le \tilde{\alpha}_j) 
= \Pr(a + c\alpha_{j-1} < a + cY_i^* \le a + c\alpha_j) 
= \Pr(\alpha_{j-1} < Y_i^* \le \alpha_j).
\]

\begin{itemize}\setlength\itemsep{0.25em}
  \item The latent component is scaled by $c$ ($\tilde{m} = c m$).
  \item The latent thresholds and covariate coefficients are shifted and scaled accordingly.
\end{itemize}

\vspace{0.4em}
\assumptionbox{Theorem 1}{
The latent systematic index $m(X_i, D_i)$ is identified only up to an arbitrary scale shift $c > 0$ and location shift $a$ (Schröder and Yitzhaki, 2017).}
\end{block}

\begin{block}{Distributional Effects (What We CAN Identify)}
Focus instead on identified probability distributions, mirroring recommendations by Hanmer and Ozan Kalkan (2013).

\vspace{0.3em}
\textbf{Target Causal Estimands:}
\[
\text{Category-specific effect: } \Delta_j = \Pr(Y_i(1) = j) - \Pr(Y_i(0) = j)
\]

\assumptionbox{Theorem 2}{
Under standard conditions (SUTVA, Conditional Ignorability $Y_i(d) \perp D_i \mid X_i$, and Positivity), the potential outcome distribution is fully identified:
\[
\Pr(Y_i(d) = j) = \mathbb{E}_X\left[ \Pr(Y_i = j \mid D_i=d, X_i) \right]
\]
  }
\assumptionbox{Proposition 1}{
  Based on Theorem 2, any causal quantity that is a functional of this
distribution including category-specific effect is also identified.
  }

\end{block}

\end{column}
\separatorcolumn

% =========================================================================
% COLUMN 2: ESTIMATION
% =========================================================================
\begin{column}{\colwidth}

\begin{block}{Standard Estimation: Parametric Ordered Probit/Logit}
\justifying
If the latent error distribution $F(\cdot)$ of $\epsilon_i$ is known to be Standard Normal or Standard Logistic, estimating the conditional probabilities is straightforward.

\begin{itemize}\setlength\itemsep{0.25em}
  \item \textbf{Ordered Probit:} Assumes $F(\cdot) = \Phi(\cdot)$ (standard normal CDF).
  \item \textbf{Ordered Logit:} Assumes $F(\cdot) = \Lambda(\cdot)$ (standard logistic CDF).
\end{itemize}

We maximize the sample log-likelihood to estimate standard parameters $\Theta = \{\Theta_m, \alpha\}$:
\[
\ell_n(\Theta) = \frac{1}{n}\sum_{i=1}^n \sum_{j=1}^J \mathbf{1}(Y_i=j)\log\left[F(\alpha_j - m(X_i, D_i)) - F(\alpha_{j-1} - m(X_i, D_i))\right]
\]
\end{block}

\begin{block}{The Core Challenge: Unknown Error Distribution}
\justifying
\begin{itemize}\setlength\itemsep{0.3em}
  \item \textbf{Unknown Base Distribution:} In real-world applications, the true latent error distribution $F(\cdot)$ is \textbf{never known} in advance.
  \item \textbf{Complex Latent Shapes:} Latent attitudes in political science can be polarized or skewed, directly violating standard symmetric and unimodal parametric shapes.
\end{itemize}

\vspace{0.4cm}
\vspace{0.4cm}
\keybox{Forcing standard normal or logistic errors when the true latent shape is complex  \\ leads to biased and unreliable causal effects $\Delta_j$.}
\end{block}

\begin{block}{The Core Idea: Normalizing Flows}
\justifying
How can we estimate the shape of $F(\cdot)$ directly from the data?

\vspace{0.3em}
\begin{itemize}\setlength\itemsep{0.2em}
    \item A \textbf{Normalizing Flow} (originally introduced by Rezende and Mohamed (2015)) is an invertible map $T_\phi$ that stretch and warp the base distribution to match the true error distribution ($\epsilon_i$):
  \[
  \epsilon_i = T_\phi(Z_i \mid D_i, X_i)
  \]
  where $\phi$ is a vector of parameters for the flow model.
  \item This gives us the cumulative distribution function (CDF) of the true error distribution analytically (Papamakarios et al., 2021):
  \[
  F(\epsilon_i \mid d, x) = \Phi\left( T_\phi^{-1}(\epsilon_{i} \mid d, x) \right)
  \]
  \item Panel (a) shows how standard ordered probit assumes an identity mapping ($T = \text{id}$), which forces the estimated latent shape to match a rigid normal distribution even when the true shape (dashed green) is bimodal.
  \item Panel (b) shows how the normalizing flow estimates $T_\phi$ directly from the data, warping the thresholds so that the density naturally "piles up" to represent multimodal true latent errors.
\end{itemize}

\vspace{0.6cm}
% {{{
\begin{center}
\begin{tikzpicture}[x=2cm,y=3.2cm,>=Latex]
  % ==================== PANEL (a): Ordered Probit (Identity) ====================
  \begin{scope}[shift={(-4.2,0)}]
    \node[above=8pt, font=\small\bfseries] at (0,4.5) {(a) Ordered probit: Assume Normal Error};
    
    % Base distribution Z
    \draw[color=PosterGray] (-2.7,3.5)--(2.7,3.5);
    \draw[thick, color=blue!80!black] plot[smooth] coordinates {
      (-2.5,3.52) (-1.8,3.62) (-1.2,3.88) (-0.6,4.18) (0,4.3) (0.6,4.18) (1.2,3.88) (1.8,3.62) (2.5,3.52)
    };
    \node[left, font=\tiny, color=PosterGray] at (-1.3,4.2) {$Z \sim N(0,1)$};
    \node[right, font=\tiny, color=PosterGray] at (1.3,4.2) {evenly spaced};
    
    % Mapping line (Middle)
    \draw[color=PosterGray] (-2.7,2.0)--(2.7,2.0);
    \node[left, font=\tiny] at (-1.5,2.3) {$\epsilon = Z$};
    \node[right, font=\tiny, color=PosterGray] at (0,2.3) {(spacing unchanged)};
    
    % Connectors & Dots
    \foreach \x in {-2.0, -1.5, -1.0, -0.5, 0.0, 0.5, 1.0, 1.5, 2.0} {
      \fill[blue!80!black] (\x,3.5) circle (3pt);
      \fill[red!80!black] (\x,2.0) circle (3pt);
      \draw[color=PosterGray!50, thin] (\x,3.5)--(\x,2.0);
    }
    
    % Target distribution (Bottom)
    \draw[color=PosterGray] (-2.7,0.5)--(2.7,0.5);
    
    % Assumed normal error shape (Red, solid)
    \draw[thick, color=red!80!black] plot[smooth] coordinates {
      (-2.5,0.52) (-1.8,0.62) (-1.2,0.88) (-0.6,1.18) (0,1.3) (0.6,1.18) (1.2,0.88) (1.8,0.62) (2.5,0.52)
    };
    \node[red!80!black,above,font=\tiny,rotate=15] at (-0.3,1.2) {assumed error};
    
    % True multimodal error shape (Green, dashed)
    \draw[thick, dashed, color=green!60!black] plot[smooth] coordinates {
      (-2.5,0.51) (-2.0,0.68) (-1.5,0.98) (-1.0,0.72) (-0.5,0.54) (0,0.51) (0.5,0.54) (1.0,0.72) (1.5,0.98) (2.0,0.68) (2.5,0.51)
    };
    \node[green!60!black,above,font=\tiny] at (-1.5,1.0) {true error};
    
    % Thresholds
    \draw[dashed, color=PosterGray] (-0.8,0.5)--(-0.8,1.1);
    \draw[dashed, color=PosterGray] (0.8,0.5)--(0.8,1.1);
    \node[below, font=\tiny] at (-0.8,0.5) {$\alpha_1 - m(d,x)$};
    \node[below, font=\tiny] at (0.8,0.5) {$\alpha_2 - m(d,x)$};
  \end{scope}

  % ==================== PANEL (b): Conditional Flow (Spline Warp) ====================
  \begin{scope}[shift={(4.2,0)}]
    \node[above=8pt, font=\small\bfseries] at (0,4.5) {(b) Conditional flow: $T_\phi$ estimated};
    
    % Base distribution Z
    \draw[color=PosterGray] (-2.7,3.5)--(2.7,3.5);
    \draw[thick, color=blue!80!black] plot[smooth] coordinates {
      (-2.5,3.52) (-1.8,3.62) (-1.2,3.88) (-0.6,4.18) (0,4.3) (0.6,4.18) (1.2,3.88) (1.8,3.62) (2.5,3.52)
    };
    \node[left, font=\tiny, color=PosterGray] at (-1.3,4.2) {$Z \sim N(0,1)$};
    \node[right, font=\tiny, color=PosterGray] at (1.3,4.2) {evenly spaced};
    
    % Mapping line (Middle)
    \draw[color=PosterGray] (-2.7,2.0)--(2.7,2.0);
    \node[left, font=\tiny] at (-1.5,2.3) {$\epsilon = T_\phi(Z \mid d, x)$};
    \node[right, font=\tiny, color=PosterGray] at (0,2.3) {(crowded $\implies$ density piles up)};
    
    % Connectors & Dots (Warped Mapping)
    \foreach \z/\e in {-2.0/-2.3, -1.5/-1.6, -1.0/-1.45, -0.5/-1.35, 0.0/0.0, 0.5/1.35, 1.0/1.45, 1.5/1.6, 2.0/2.3} {
      \fill[blue!80!black] (\z,3.5) circle (3pt);
      \fill[red!80!black] (\e,2.0) circle (3pt);
      \draw[color=PosterGray!50, thin] (\z,3.5)--(\e,2.0);
    }
    
    % Target distribution (Bottom Bimodal)
    \draw[color=PosterGray] (-2.7,0.5)--(2.7,0.5);
    
    % Estimated error shape (Red, solid - matches the true shape)
    \draw[thick, color=red!80!black] plot[smooth] coordinates {
      (-2.5,0.51) (-2.0,0.68) (-1.5,0.98) (-1.0,0.72) (-0.5,0.54) (0,0.51) (0.5,0.54) (1.0,0.72) (1.5,0.98) (2.0,0.68) (2.5,0.51)
    };
    \node[red!80!black,above,font=\tiny] at (1.5,1.0) {estimated error};
    
    % True multimodal error shape (Green, dashed) - slightly shifted/separated
    \draw[thick, dashed, color=green!60!black] plot[smooth] coordinates {
      (-2.5,0.51) (-2.0,0.71) (-1.5,1.03) (-1.0,0.69) (-0.5,0.52) (0,0.49) (0.5,0.52) (1.0,0.69) (1.5,1.03) (2.0,0.71) (2.5,0.51)
    };
    \node[green!60!black,below,font=\tiny] at (-1.5,1.3) {true error};
    
    % Thresholds
    \draw[dashed, color=PosterGray] (-0.8,0.5)--(-0.8,0.75);
    \draw[dashed, color=PosterGray] (0.8,0.5)--(0.8,0.75);
    \node[below, font=\tiny] at (-0.8,0.5) {$\alpha_1 - m(d,x)$};
    \node[below, font=\tiny] at (0.8,0.5) {$\alpha_2 - m(d,x)$};
  \end{scope}
\end{tikzpicture}
\end{center}
% }}}
\end{block}

\begin{block}{Likelihood and Consistency}
\justifying
  Substituting the flow-based CDF, the parameters $\Theta = \{\Theta_m, \phi, \alpha\}$ maximize the sample log-likelihood for $m = m(X_{i}, D_{i})$:

\begin{multline*}
\ell_n(\Theta) = \frac{1}{n}\sum_{i=1}^n \sum_{j=1}^J \mathbf{1}(Y_i = j) \\ \log \left[ \Phi\left(T_\phi^{-1}(\alpha_j - m \mid D_i, X_i)\right) - \Phi\left(T_\phi^{-1}(\alpha_{j-1} - m \mid D_i, X_i)\right) \right]
\end{multline*}

\vspace{0.4cm}
\textbf{Assumption C.3 (Flow Approximation \& Identification):}
The family of distributions that can be approximated by the flow is rich enough to approximate the true error CDF $F_0$ arbitrarily well. 

\vspace{0.3em}
\textbf{Assumption C.4 (Regularity Conditions):}
The parameter space is compact, the likelihood satisfies standard continuity and differentiability conditions.

\vspace{0.5em}
\end{block}

\end{column}
\separatorcolumn

% =========================================================================
% COLUMN 3: EVIDENCE
% =========================================================================
\begin{column}{\colwidth}

\assumptionbox{Theorem C.2 (Consistency of the Flow Estimator)}{
Suppose the latent outcome model and threshold structure are correctly specified, and Assumptions C.3--C.4 hold. Then, as $n \to \infty$, the maximum likelihood estimator is consistent up to a positive affine transformation:
\[
\| \widehat{\Theta} - \Theta_0 \| = o_p(1).
\]
}

\begin{itemize}
  \item Therefore, the flow-based estimators consistently estimate $\Delta_{j}$ 
\end{itemize}

\begin{block}{Why Monotone Rational-Quadratic Splines (RQS)?}
\justifying
\begin{itemize}\setlength\itemsep{0.3em}
  \item Many of normalizing flows, such as those based on affine coupling (Dinh et al., 2014) or RealNVP (Dinh et al., 2017), are designed primarily to output latent density functions $f(\epsilon)$. 
  \item Evaluating the threshold probabilities in ordered models requires the cumulative distribution function (CDF) $F(\epsilon) = \int_{-\infty}^{\epsilon} f(t) dt$, which is computationally intractable under density-only flows without expensive numerical integration.
  \item Monotone RQS (Durkan et al., 2019) parameterizes the transformation $T_\phi$ using continuously differentiable, strictly increasing piecewise rational-quadratic splines, yielding an \textbf{analytical, closed-form inverse $T_\phi^{-1}$ and exact CDF}.
  \item This mathematical structure allows us to evaluate the ordered likelihood exactly, enabling fast, stable, and exact maximum-likelihood estimation.
\end{itemize}
\end{block}

\begin{block}{Nonlinear Latent Model Extension}
\justifying
When the systematic relationship $m(X_{i}, D_{i})$ is highly non-linear or interactive:
\begin{itemize}\setlength\itemsep{0.3em}
  \item The normalizing flow $T_\phi$ can condition directly on a nonparametric representation of $(D_i, X_i)$ without altering the likelihood setup.
  \item This ``model-free'' specification bypasses parametric linear modeling entirely, naturally capturing arbitrary treatment heterogeneities and complex covariate interactions.
\end{itemize}
\end{block}

\begin{block}{Monte Carlo Simulation}

\begin{itemize}\setlength\itemsep{0.25em}
  \item \textbf{Polarization:} a bimodal latent error distribution.
  \item \textbf{Nonlinearity:} nonlinear treatment-response structure.
\end{itemize}

\begin{figure}
  \centering
  \maybeincludegraphics{width=0.98\textwidth}{../figures/fig_polarized_mixture_N2000.pdf}
  \caption{Category-specific effects under a polarized latent distribution ($n = 2000$).}
\end{figure}
\vspace{0.15em}
\begin{figure}
  \centering
  \maybeincludegraphics{width=0.98\textwidth}{../figures/fig_app_nonlinear_moderates_N2000.pdf}
  \caption{Category-specific effects under a nonlinear structure ($n = 2000$).}
\end{figure}
\vspace{-0.2em}
\begin{itemize}\setlength\itemsep{0.22em}
  \item Parametric error misspecification distorts category-level effects.
  \item Structured flows improve robustness with little efficiency loss.
  \item Model-free flows are most flexible but noisier in finite samples.
\end{itemize}
\end{block}

\begin{block}{Conclusion}
\begin{enumerate}\setlength\itemsep{0.35em}
  \item \textbf{Cardinalized mean effects are not identified with ordinal outcomes.}
  \item \textbf{Category-specific effects are identified under standard identification assumptions.}
\end{enumerate}
\end{block}

\begin{block}{Contact and Materials}
\textbf{Chanhyuk Park}\par
Washington University in St. Louis\par
Email: \texttt{chanhyuk@wustl.edu}\par
\vspace{0.2em}
Code \& Data: \url{https://github.com/chanhyuk58/ordinal_flow}
\end{block}

\end{column}
\separatorcolumn
\end{columns}
\end{frame}
\end{document}
